\chapter{Рабочая неделя и доказательство}

\section{Бюджет 12--15 часов}
Базовая неделя занимает 12 часов.
\begin{center}
\begin{tabularx}{\textwidth}{@{}P{0.38\textwidth}P{0.1\textwidth}Y@{}}
\toprule
Работа & Время & Результат\\
\midrule
Основная лемма: три сессии & 6 ч & Письменные попытки, доказанный переход
или контрпример\\
Литература по текущему препятствию & 2 ч & Точная внешняя теорема,
проверка гипотез и сравнение\\
Чистовой математический текст & 2 ч & Исправленная формулировка
и связный фрагмент доказательства\\
Примеры, вычисления или Lean & 1 ч & Воспроизводимая проверка
одного спорного места\\
Обсуждение и журнал & 1 ч & Одно главное препятствие и задача
следующей недели\\
\bottomrule
\end{tabularx}
\end{center}
При 15 часах добавить два часа к доказательству и час к проверкам
или внешнему обсуждению. В первые две недели действует бюджет отбора:
4 часа --- близкие теоремы и доказательства, 3 часа --- попытки сведения,
3 часа --- формулировка и сравнение, 1 час --- примеры,
1 час --- обсуждение и журнал. После обоснования дополнительного вопроса
применяется основная таблица. Семинар и переписка
с исследователем входят в общий бюджет.

Параллельных обязательных курсов в этом цикле нет. Чтение Любича,
Михайлова--Пасси, Хэтчера и Милнора оплачивается из двух часов литературы.
Увеличение этого бюджета на конкретной неделе записывается вместе
с названием закрываемой зависимости.

\section{Одна исследовательская сессия}
\begin{enumerate}
\item Назвать один переход, который должен быть получен за сессию.
\item Выписать все доступные гипотезы и проверить простейший пример.
\item Работать над аргументом 60--90 минут. При остановке записать
      последнее корректное утверждение и точный пробел.
\item Последние 15--20 минут посвятить тексту: доказанный шаг, неудачная
      оценка, контрпример или уточнение постановки.
\end{enumerate}
Если найдено неверное утверждение, сохранить его вместе с контрпримером.
Рабочие файлы первой темы перечислены в \repo{research/README.md}.

\section{Разбор нужного известного доказательства}
Сначала выделить используемую лемму и попытаться доказать нужный частный
случай. Затем прочитать источник, проверить условия и закрыть книгу.
Восстановить аргумент письменно; через неделю повторить ключевой переход
и проверить изменённый пример.

Для внешнего результата сохраняются автор, версия, номер теоремы,
страница и проверка гипотез. Полное воспроизведение крупной теории
назначается отдельной задачей только при необходимости для основного
доказательства.

\section{Состояния работы}
В журнале различать: \emph{кандидат}, \emph{в работе},
\emph{доказано на бумаге}, \emph{проверено извне},
\emph{опровергнуто}, \emph{сведено к известному}.
Новизна фиксируется отдельно: \emph{не проверена},
\emph{найден предшественник}, \emph{элементарное следствие},
\emph{добавление предварительно обосновано},
\emph{сравнение обсуждено со специалистом}.
Научный вес описывается отдельным абзацем с примером предполагаемой пользы.
Полученный отзыв записывается дословно по существу;
отсутствие ответа остаётся отсутствием проверки.

\section{ИИ и Lean}
Использовать ИИ для одной конкретной леммы, поиска ошибки, проверки
кванторов или подготовки формализации. Предложенный аргумент получает
статус после проверки каждого перехода и используемых источников.
Автора идеи, текст доказательства и способ проверки отмечать в карточке.

Формализовать устойчивое бумажное утверждение. В базовом бюджете
на Lean и вычисления приходится один час в неделю; при 15 часах
можно выделить до двух. После формализации сверить определения,
проверить зависимости через \repo{#print axioms} и собрать модуль.
Техническая задержка переносится в отдельный список.

\section{Правило остановки}
После двух сессий без продвижения выписать препятствие как отдельную
лемму. После четырёх --- проверить модельный случай, источник или получить
предметный разбор. После двух недель с тем же пробелом уменьшить общность
или изменить конструкцию. Решение записать с математической причиной.
